\subsection{A compact matrix representative for every unitary orbital vector}
\label{actual-unitary-density}

The remaining density transfer requires a coefficient for all relevant
heads at once.  A section of the reduced-matrix map, followed by
integration over its actual fibres, supplies it.  The positive Gram
condition is essential: it makes the complement of the two flag frames
positive definite and hence makes the fibre integration proper.

Fix the original real-place data, with $k\ge1$, $r\ge0$, $n=r+1$,
negative middle vectors $z,z'$, and
$T=-h(z,z)>0$, $T'=-h(z',z')>0$.  The Hermitian form is linear in
its first variable.  Write $f=(e_1,\ldots,e_r,z)$,
$f'=(e_1,\ldots,e_r,z')$ and
\[
 Z(g)_{ab}=h(gf'_b,f_a),\qquad
 G(Z)=|\det Z|^2-TT'|\det Z_{[r],[r]}|^2.
\]
All groups, compact representations and Haar measures below are the
prescribed ones.  In particular the compact Haar measures are not
renormalized to have mass one.  Let $N_0(z)$ be the kernel of the
flag-action map $u_z:N\to N_n$; its elements fix every $f_a$.
Let $K_z$ be the original compact group fixing the hyperbolic pairs
and $z$.  The fixed base frame consists of a positive unit vector
$\nu'\perp z'$ in the middle space and an isometry $m_0$ fixing
the hyperbolic pairs, with $m_0z'=\sqrt{T'/T}\,z$.

\begin{theorem}\label{unitary-density-compact-representative}
For every finite-dimensional coefficient space
$V=\operatorname{Hom}(W',W)$ and every original compact coefficient
$F$, there is $\Phi_F\in C_c^\infty(\operatorname{Mat}_n(\mathbb C),V)$
whose support consists of invertible matrices and for which
\begin{equation}\label{unitary-density-orbital-identity}
       V_F(a)=\operatorname{PhiTop}_{\beta_U}(a;\Phi_F)
\end{equation}
for every relevant block head $a$ of total size $n$.
The character is the original $\beta_U=2i$, and the right-hand side
uses the original planar Lebesgue measure in its restricted lower
and full upper root coordinates.  No internal leading minor is
assumed nonzero.  The assertion includes $r=0$ and $k=1$.

Here ``original compact coefficient'' means that $F$ is smooth on
the unitary group and there are $R<\infty$ and $\epsilon>0$ such
that $F(g)\ne0$ implies $\|g\|\le R$ and $G(Z(g))\ge\epsilon$.
The construction preserves the original representation factors,
quotient stabilizer measures and the two positive Haar constants.
\end{theorem}

\begin{proof}
We give the section, its properness, and the exact orbital comparison.

\paragraph{The global section.}
Work first in the space spanned by
$(e_i,e_i^*,z',\nu')$, with coordinates $(a,b,\alpha,\beta)$ and
form
\[
 h((a,b,\alpha,\beta),(a',b',\alpha',\beta'))
 =a^t\overline{b'}+b^t\overline{a'}
       -T'\alpha\overline{\alpha'}+\beta\overline{\beta'}.
\]
For $G(Z)>0$, put
\[
 \kappa=\operatorname{adj}(Z)e_n,\quad N=\kappa^*\kappa,
 \quad\varphi=\frac{\det Z}{|\det Z|},\quad
 s=\sqrt{G(Z)/T},\quad s'=\sqrt{G(Z)/T'}.
\]
Thus $\det Z\ne0$, $N>0$ and $s,s'>0$.  Define
\[
 \alpha_b=-\frac{Z_{nb}}{\sqrt{TT'}},\qquad
 \beta_b=\frac{\overline{\kappa_b}}N s\varphi,
 \qquad A'=-T'e_ne_n^*.
\]
Let $R$ be the usual conjugate-first Gram defect of the middle
vectors $(\alpha_b,\beta_b)$ from $A'$; explicitly
\[
 R_{cb}=A'_{cb}
       +T'\overline{\alpha_c}\alpha_b
       -\overline{\beta_c}\beta_b.
\]
Set
\[
 Q=\frac{\kappa\kappa^*}{N},\quad P=I-Q,\quad
 K=\tfrac12PRP+PRQ,\quad M=(Z^{-1})^*K.
\]
The identities $Z\kappa=(\det Z)e_n$ and
$\kappa_n=\det Z_{[r],[r]}$ give
\[
 \kappa^*R\kappa
   =-T'|\kappa_n|^2+|\det Z|^2/T-G(Z)/T=0.
\]
Consequently $QRQ=0$, $K+K^*=R$, $QK=0$, and the last row of
$M$ is zero.  The vectors
\[
 x_b=(M_{[r],b},Z_{[r],b},\alpha_b,\beta_b)
\]
therefore have Gram matrix $A'$ and satisfy
$h(x_b,f_a^0)=Z_{ab}$, where
$f_a^0=e_a$ for $a<n$ and $f_n^0=\sqrt{T/T'}\,z'$.
This verifies the required reduced matrix without dividing by
any internal minor.

Complete these vectors by the following explicit indefinite
Gram--Schmidt procedure.  For $i<n$ put
\[
 y_i^0=\sum_a((Z^{-1})^*)_{ai}f_a^0,\qquad
 y_i=y_i^0-\tfrac12\sum_{j<n}h(y_i^0,y_j^0)x_j.
\]
Then $h(y_i,x_j)=\delta_{ij}$ for $j<n$,
$h(y_i,x_n)=0$, and $h(y_i,y_j)=0$.
Put $c=\sum_a\overline{\operatorname{adj}(Z)_{na}}f_a^0$ and
\[
 p=c-\sum_{i<n}h(c,y_i)x_i-\sum_{i<n}h(c,x_i)y_i
                       +\frac{h(c,x_n)}{T'}x_n,
 \qquad v=\frac{\varphi}{s'}p.
\]
The adjugate identity gives $h(c,x_b)=\overline{\det Z}\,
\delta_{nb}$ and $h(c,c)=-T|\kappa_n|^2$.  Thus $p$ is
orthogonal to the $x_b,y_i$, its norm is $G(Z)/T'$, and $v$ is
a positive unit vector.  There is now a unique isometry $\sigma_Z$
of this $2n$-dimensional space sending
\[
       (e_i,e_i^*,z',\nu')\longmapsto(x_i,y_i,x_n,v).
\]
Extend it by the identity on the middle space
$W_0=\{z',\nu'\}^{\perp}$, and set $\gamma_Z=m_0\sigma_Z$.
Every displayed formula is smooth on $G>0$.

We check the prescribed frame normalization, which is needed for
vector-valued orbital integrals.  The right cofactor vector is
$\sum_b\kappa_b\gamma_Z f'_b$; its middle projection orthogonal
to $z$ is $s\varphi\,m_0\nu'$.  The corresponding inverse left
cofactor has projection $(s'/\varphi)\nu'$.  The two positive
normalizing scalars are therefore $s/|\det Z|$ and
$s'/|\det Z|$, exactly as required by the original $D$-frame.
For the third condition, the Gram matrix of $(f^0,x)$ in the
conjugate-first convention is
\begin{equation}\label{unitary-density-paired-gram}
       \begin{pmatrix}-Te_ne_n^*&Z\\Z^*&-T'e_ne_n^*\end{pmatrix},
       \qquad \det=(-1)^nG(Z).
\end{equation}
The two frames thus span the whole $2n$-dimensional space above.
Their common orthogonal complement is $W_0$, on which $\sigma_Z$
is the identity; the third frame condition follows.  Also
$\sigma_Z$ commutes with every isometry fixing the hyperbolic
pairs, $z'$ and $\nu'$, because such an isometry is the identity
on that $2n$-dimensional space.  This will remove compact-frame
ambiguity below.

\paragraph{Properness and smooth integration.}
For arbitrary $g$ with $G(Z(g))>0$, the matrix in
\eqref{unitary-density-paired-gram} is nondegenerate and has
signature $(n,n)$.  Indeed, scale both diagonal blocks by
$t\in[0,1]$; its determinant remains
$(-1)^n(|\det Z|^2-t^2TT'|\kappa_n|^2)\ne0$, and at $t=0$
the signature is $(n,n)$.  The paired frame $(f,gf')$ therefore
has a positive definite orthogonal complement of dimension $k-1$.

Let $g$ vary in a compact subset of $G(Z(g))>0$.  If
\[
       y=h^{-1}n_0^{-1}g n'_0h'
\]
lies in the fixed compact support of $F$, where
$h\in K_z$, $h'\in K_{z'}$, $n_0\in N_0(z)$,
$n'_0\in N_0(z')$, then all four integration variables lie in
one compact set.  To see this, $t=h^{-1}n_0^{-1}$ fixes $f$
and sends $gf'$ to $yf'$.  The two paired Gram matrices agree
and have inverses bounded uniformly on these compact sets.
Thus $t$ is bounded on their spans.  On their positive
orthogonal complements it is an isometry; the positive forms
and Euclidean norms are uniformly equivalent, so $t$ is
bounded there too.  The orthogonal projections are uniformly
bounded because the paired Gram inverses are.  Hence $t$, then
$n_0$, are bounded.  Apply the same argument to $g^{-1},y^{-1}$
to bound $n'_0$.  The groups $N_0(z),N_0(z')$ are closed, so the
bounded sets in question have compact closure.

Choose arbitrary nonzero Haar measures $\mu_0,\mu'_0$ on these
two kernels and use the original compact Haar measures.  Define
\begin{equation}\label{unitary-density-fibre-integral}
 I_F(g)=\int
 \pi(h)F(h^{-1}n_0^{-1}g n'_0h')\pi'(h'^{-1})
           \,d\mu_W(h)\,d\mu_0(n_0)
                 \,d\mu_{W'}(h')\,d\mu'_0(n'_0).
\end{equation}
The preceding compactness is local and uniform in $g$.  It
justifies every differentiation under this integral, using
local smooth extensions of $F$ off the unitary group.  Only
continuity of the two compact representations is needed: they
are bounded and the derivatives act on $F$ and $\gamma_Z$.
Thus $I_F(\gamma_Z)$ is smooth on $G>0$.

The fibre preserves $Z$.  Nonvanishing of this integral implies
that some $y$ in the compact support of $F$ has $Z(y)=Z$.
Consequently its support is contained in a compact set of
matrices satisfying $G(Z)\ge\epsilon$.  Extending it by zero
on $G\le0$ is smooth and compactly supported, with
$|\det Z|^2\ge\epsilon$ on its support.

\paragraph{Why the section has the original frame.}
We record the precise uniqueness needed to compare the fibre
integral with the original orbital vector.  If $g_1,g_2$ have
the same reduced matrix with positive Gram, the isometry
$(f,g_1f')\mapsto(f,g_2f')$ extends to the whole space: extend
it by any isometry between their positive orthogonal
complements.  It follows that
\[
       g_2=(h_A n_A)^{-1}g_1(h_B n_B),
       \qquad h_A\in K_z,\ h_B\in K_{z'},\
                         n_A\in N_0(z),\ n_B\in N_0(z').
\]
Here the decomposition of the pointwise flag stabilizer is
obtained by its induced isometry on the positive middle
quotient; removing that isometry leaves the indicated
unipotent kernel.

Suppose $g_1=m_0\sigma$ is a $D$-frame, with $\sigma$ fixing
$W_0$ and commuting with the compact group fixing $\nu'$, as
our constructed section does, and $g_2$ is also a $D$-frame.
The first positive cofactor normalization forces
$h_A^{-1}m_0\nu'=m_0\nu'$: its possible scalar is positive
real and both vectors have norm one.  The second similarly
forces $h_B^{-1}\nu'=\nu'$.  For $w\in W_0$, apply the third
frame condition to $(h_A n_A)^{-1}m_0w$; it gives
\[
       h_A^{-1}m_0w=m_0h_B^{-1}w.
\]
In making this projection the $N_0$ terms disappear, since
they act trivially on the middle quotient of the orthogonal
complement of the isotropic flag.  The same identity holds on
the hyperbolic pairs and on $z',\nu'$, which these compact
elements already fix.  Hence $h_A^{-1}m_0=m_0h_B^{-1}$ and
$h_A^{-1}g_1h_B=g_1$.  It follows that
$g_2=n_A^{-1}g_1n_B$ and therefore $I_F(g_2)=I_F(g_1)$ by
left invariance of $\mu_0,\mu'_0$.

The same conclusion holds for
$g_1=\tau_z(v)^{-1}m_0\sigma\tau_{z'}(u)$, where $v,u\in N_n$
and $\tau_z$ is the original upper-flag section.  Indeed,
$\tau_z(v)$ commutes with $K_z$, normalizes $N_0(z)$, and
preserves the projection onto the middle complement of $z$.
Its action on the flag is $f\mapsto fv$.  The adjugate
identities consequently show that the $D$-frame conditions
are preserved by left and right $\tau$ translations.  Undo
the two translations, apply the preceding compact cancellation,
and redo them.  This proves the equality of the two fibre
integrals with no representation scalar or conjugation factor.
These are direct identities for the prescribed Eichler
section, not a choice of a different flag action.

\paragraph{The measures and the character.}
The map $(u,n_0)\mapsto\tau_z(u)n_0$ is a smooth
coordinate decomposition of $N$, with quotient $N_n$.
The product of upper-entry planar Lebesgue measure and
$\mu_0$ pushes forward to a Haar measure: left translation
acts by left translation on $N_n$ and by a left translation
on each $N_0$ fibre.  This follows directly by applying
$u_z$ and then solving for the kernel coordinate.  The
coordinate map and its inverse are continuous, so the
pushforward is locally finite and nonzero.  Haar uniqueness
gives the exact positive constants $c_L,c_R$ satisfying
\begin{equation}\label{unitary-density-haar-constants}
 d\mu_N=c_L\,d u\,d\mu_0,
 \qquad d\mu_N=c_R\,d u'\,d\mu'_0.
\end{equation}
No value for either constant is presumed.  Changing a kernel
Haar measure changes its constant inversely and leaves the
construction unchanged.

At a relevant head matrix $H_a$, decompose the two upper
groups by their stabilizer roots.  The complementary lower
roots are precisely those between the reversed blocks of
$a$, and all upper roots remain.  In increasing root-height
coordinates this decomposition is triangular with diagonal
entries one.  Thus it preserves the product planar measure.
The original stabilizer measure is the pushforward of these
same discarded root coordinates.  The normalized quotient
cut has integral one along them, so that integration removes
exactly that factor.  Under an $N_0$ change of representative,
the flag-action map on the stabilizer is unchanged; it is
injective there because a unipotent element in the compact
joint-frame stabilizer is the identity.  Hence its specified
measure also has no scalar ambiguity.  Its image is the full
discarded root group: for $v^*H_a u=H_a$, the actions $v,u$ on
the paired frames define an isometry, which extends by the
identity on their positive complement.  It is unipotent, so
its compact middle quotients are the identity and it belongs
to $N\cap\gamma N\gamma^{-1}$.  Thus the prescribed inverse
flag-action map gives exactly this root measure, not a
Haar measure merely proportional to it.

All changes of order are absolutely convergent.  The
restricted block-lower/full-upper coordinates are bounded
when $LH_aU$ is in a compact subset of invertible matrices,
by block Gaussian elimination.  The fibre properness just
proved bounds the remaining unipotent coordinates, while
the compact groups have finite Haar measure.  This also
constructs an admissible normalized quotient cut in the
discarded root coordinates.

Since $\ell_z(\tau_z(v))$ is the superdiagonal sum and
$\ell_z(N_0(z))=0$, the left character after $L=v^*$ is
$\psi_{c_{\rm Std}}(\overline{\theta_-(L)})
=\psi_{\beta_U}(\theta_-(L))$; the inverse right character
is $\overline{\psi_{c_{\rm Std}}(\theta(U))}
=\psi_{\beta_U}(\theta(U))$.
Here $c_{\rm Std}=-2i$ and $\beta_U=-c_{\rm Std}=2i$.
The quotient integral is therefore exactly
\[
 c_Lc_R\int_{L,U}
       \psi_{\beta_U}(\theta_-(L)+\theta(U))
              I_F(\gamma_{LH_aU})\,dL\,dU.
\]
The preceding frame cancellation identifies this integrand
with the one obtained from the original chosen $D$-frame.
The formula proves \eqref{unitary-density-orbital-identity}
with $\Phi_F=c_Lc_R I_F(\gamma_Z)$ on $G>0$, and zero
elsewhere.  If $G(H_a)\le0$, both sides vanish: $G(LH_aU)=G(H_a)$.
This includes every relevant head and completes the construction.
\end{proof}

\begin{corollary}\label{unitary-density-every-head}
For every original compact coefficient $F$, if its regular
radial transform is zero, then $V_F(a)=0$ at every relevant
head $a$.
\end{corollary}

\begin{proof}
The established big-cell pointwise theorem makes every regular
orbital integral of $\Phi_F$ vanish.  The compact complex density
theorem of Section~\ref{actual-terminal-density} then makes all
its relevant orbital integrals vanish.  Apply
\eqref{unitary-density-orbital-identity}.  This is a pointwise
statement, including the lower-dimensional heads sampled
arithmetically; no null-set replacement is used.
Heads failing the total-size, positive-block-size or nonzero-entry
conditions have $V_F(a)=0$ by its original definition, so the
universal quantifier over all heads in the source density contract
is also covered.
\end{proof}

\paragraph{Scope.}
The theorem supplies the original compact matrix-representative
contract and the resulting every-head density contract
\cite{unitary-density-source}.  Its local compactness constants
may depend on the place data, rank and compact support.  Only
an exact zero conclusion is transferred, so none enters a
uniform arithmetic norm bound.  Completing the terminal
operator transfer still requires the compact top-lift
construction with the prescribed operator and all word
parameters.  This result alone does not prove the full kernel
estimate or complete W06/W08.
